\section{Tesztelés és validáció}

Az alkalmazás tesztelése elsősorban funkcionális szempontból történt.
A tesztelési fázis célja annak ellenőrzése volt, hogy a felhasználói folyamatok
hibamentesen végigvihetők-e, valamint hogy a méréshez kapcsolódó adatok a várt
formátumban menthetők és tölthetők vissza.

\subsection{Fizikai és emulált tesztkörnyezet}

A tesztelés két fő környezetben zajlott. A fejlesztési fázis nagy részében az
Android Virtual Device (AVD) emulátorok használata tette lehetővé a gyors iterációt
\cite{androidemulator}. Ennek segítségével a felhasználói felület (UI) vizuális elemei,
a képernyőméretekhez való alkalmazkodás és a fájlkezelési funkciók anélkül
voltak tesztelhetők, hogy minden alkalommal fizikai teszteszközt kellett
volna csatlakoztatni.

A tesztelés megtörtén fizikai eszközön is. Erre azért volt szükség, hogy a
Bluetooth Low Energy (BLE) kapcsolat stabilitása, az operációs rendszer szintű
jogosultságkérések, a hardveres kameraelőnézet, valamint a valós idejű,
aszinkron szenzoradatok fogadása valós körülmények között is igazolható legyen.

\subsection{Adatgenerálás (Mocking)}

Mivel a zokni nem volt elérhető a tesztelés minden fázisában, ezért az alkalmazás egy
beépítettem egy tesztadat-generátort (\texttt{MockDataGenerator}) , amely képes a
fizikai szenzor adatáramlását szimulálni \cite{syntheticdata}.

A generátor nem csupán statikus értékeket állít elő, hanem a kiválasztott rehabilitációs
feladat (pl. helyben járás, tandem állás) típusához igazodó, matematikailag modellezett
jelsorozatokat generál:
\begin{itemize}
\item Dinamikus mozgás szimulálása:  Helyben járás esetén az algoritmus egy
belső időzítő alapján végigmegy a lépésciklus fázisain.
Különálló gyorsulás- és nyomásértékeket rendel a sarok földetéréséhez,
a talp gördüléséhez, az elrugaszkodáshoz és a lengési fázishoz.
\item Súlypontvándorlás: Statikus feladatoknál a generátor trigonometrikus függvények
(\texttt{sin}) segítségével szimulálja a test folyamatos, hullámszerű előre-hátra dőlését.
\item Zaj és patológiás tünetek: A valósághűség érdekében a rendszer random zajt generál
az adatokba a szenzorok 10-bites hardveres limitjein belül. Emellett dedikált logikát
tartalmaz a Parkinson-kórra jellemző mikro-remegések és ritka, véletlenszerű lefagyásos
(Freezing of Gait) epizódok szimulálására is.
\end{itemize}

A generált adatok használata lehetővé tette, hogy az adatgyűjtés folyamata, az UI
diagramok frissülése és az elemzési logika (pl. lépésdetektálás, varianciaszámítás)
folyamatosan tesztelhető maradjon. Mivel a szintetikus szenzoradatok felépítése
meghatározott, az adatfeldolgozó algoritmusok pontossága és helyessége a tesztelés
során ellenőrizhető volt.

